\section{What the benchmark changed in the skill}
\label{sec:fix}

The benchmark was built to grade the skill, and it ended up debugging it twice. Each version below is a git tag
in the repository, and each was run on the full task set.

\paragraph{v1: the loop stopped after one win.} In the pilot and the first batch, skill v1 agents on
\texttt{logstats} made one large rewrite (a 400$\times$ speedup is available in a single edit), saw a big gain,
and stopped. Their own words were close to ``the loop stopped because the gain was large''. The skill's patience
rule (stop after $k$ rounds without a win) was ignored, and so was its bookkeeping: across the six
\texttt{logstats} runs the results log held 0 to 2 rounds. The scores were
fine. The process the skill exists to enforce had not happened.

\paragraph{v2: one command per round, and a big win is not a stop rule.} We moved the bookkeeping out of the
agent's hands. \texttt{hc.py setup} freezes the score command, regex, direction, gates, frozen paths, budget and
patience in a config file; \texttt{hc.py baseline} logs round 0; and \texttt{hc.py round "<change>"} does a whole
round in one call: commit, run gates, measure, compare against the current best, keep or \texttt{git revert},
append a log row, and print \texttt{NEXT} or \texttt{STOP}. The skill text says outright that a big win is not a
stop rule and adds a section on overfitting the evaluation set. With v2, Sonnet runs logged a median of 7 rounds
(2 to 20) and Opus runs 13.5. On the three original tasks, v2's scores stayed level with v1 and the controls
(Table~\ref{tab:sonnet}). On Opus, v2 was the arm that did not weaken the test suite (Section~\ref{sec:results}).

\paragraph{v3: the noisy task found two bugs in the helper.} v2 did worst of all arms on \texttt{noisytune}
(1.31$\times$), and two of its six runs ended exactly at the baseline settings. Reading the transcripts and
replaying rounds on a scratch copy found two causes, both in \texttt{hc.py}, not in the agents.

\begin{enumerate}
\item \emph{Stale bytecode.} Agents chained \texttt{round} calls quickly. A revert and the next edit to
  \texttt{config.py} could land in the same second with the same file size. Python's bytecode cache is keyed on
  the source's modification time in whole seconds and its size, so the next measurement imported the previous
  round's settings. A real change was scored as the old configuration and read as noise. v3's \texttt{round}
  deletes untracked \texttt{\_\_pycache\_\_} folders before measuring and waits out the current second after a
  revert.
\item \emph{A noise band that ignored the run count.} v2 called a change a win only if its gain cleared
  $2(\mathrm{MAD}_{\text{best}} + \mathrm{MAD}_{\text{cand}})$ relative to the best. The median absolute deviation
  does not shrink as runs are added, so no number of runs could resolve a gain smaller than the band. At 8\% noise
  the band was 15 to 27\%. v3 uses two standard errors of the
  difference of medians, estimating each median's standard error as $1.2533 \cdot 1.4826\,\mathrm{MAD}/\sqrt{n}$,
  so the band shrinks with $\sqrt{n}$. \texttt{hc.py baseline} now prints the band at the configured run count
  and, when it is wider than the minimum gain, a \texttt{NOISY} line with the number of runs that would resolve
  it.
\end{enumerate}

On a scratch copy of the task under v3, raising one cache setting read as a 27.6\% win, and a 2\% real gain
correctly read as noise at the measured band. On the benchmark, v3 lifted \texttt{noisytune} from
1.31$\times$ to 1.43$\times$ with no run left at baseline, held the other tasks level (395$\times$, $+$22.7
points, 2.88$\times$), and logged a median of 6 rounds per run. It did not close the gap to the free-form
\textsf{loop} arm (1.65$\times$), for the reasons in Section~\ref{sec:results}.

Both bugs are the kind a skill author would not find by reading the skill. They surface only when many agents
use the helper in ways its author did not, on a task where they matter. That is an argument for benchmarking
process instructions the way we benchmark models.
